\documentclass[11pt,a4paper]{article}
\usepackage{DDTA_METHODOLOGY_GUIDE_STYLE_R1}
\usepackage{paracol}

\hypersetup{
  pdftitle={DDTA DermaTriage - Documentation and Base Analysis Case Study},
  pdfauthor={Documentation-Driven Threat Analysis research project}
}

% -----------------------------------------------------------------------------
% PAGE-INTEGRITY HASHES
% Generated over the canonical LaTeX payload between PAGE-CONTENT markers.
% Hash lines themselves are excluded. On the integrity-index page, hashes of
% preceding pages are resolved before hashing, while the page's own hash is
% represented by the canonical token <SELF> to avoid self-reference.
% -----------------------------------------------------------------------------
\newcommand{\PageMDfive}[1]{\csname PageMDfive#1\endcsname}
%<PAGE-HASH-DEFS>
\expandafter\def\csname PageMDfive1\endcsname{ccabe4bed7cca4d2f68a2005e42f36e3}
\expandafter\def\csname PageMDfive2\endcsname{0fb76dc0491da2793823b20571e5268b}
\expandafter\def\csname PageMDfive3\endcsname{6d31dd026458ea1ce9ff658056f4005d}
\expandafter\def\csname PageMDfive4\endcsname{875bd4afdf8e6be8b1995cc234253f2e}
\expandafter\def\csname PageMDfive5\endcsname{cb96dc22bf50d9f157f55b0374c9ca91}
\expandafter\def\csname PageMDfive6\endcsname{eec204cc9e490ffbdb154dd3ad79505c}
\expandafter\def\csname PageMDfive7\endcsname{26ab3bc6525c4c4c7ac958045a2487d9}
%</PAGE-HASH-DEFS>

\newcommand{\PageHashFooter}[1]{%
  \fancyfoot[R]{\sffamily\scriptsize\color{DDTAGray}MD5 contenuto: \texttt{#1}}%
}

\newcommand{\IntegrityRow}[3]{#1 & #2 & {\footnotesize\ttfamily #3} \\
}

\begin{document}
\selectlanguage{italian}
\fancyhead[L]{\sffamily\small\color{DDTAGray}DDTA - DermaTriage Case Study}
\fancyhead[R]{\sffamily\small\color{DDTABlue}Documentation + Base Analysis}
\fancyfoot[L]{\sffamily\scriptsize\color{DDTAGray}Documentation 2/3 - Base Analysis 1/3}

%<PAGE-DESC:1>Frontespizio e struttura del case study
%<PAGE-CONTENT:1>
\begin{titlepage}
\thispagestyle{empty}
\vspace*{10mm}
{\sffamily\bfseries\color{DDTANavy}\fontsize{15}{18}\selectfont Documentation-Driven Threat Analysis}\par
\vspace{5mm}
{\sffamily\bfseries\color{DDTANavy}\fontsize{29}{33}\selectfont DermaTriage Case Study}\par
\vspace{2mm}
{\sffamily\color{DDTABlue}\fontsize{18}{22}\selectfont Documentation + Base Analysis}\par

\vspace{12mm}
\begin{tcolorbox}[enhanced,arc=3pt,boxrule=0pt,colback=DDTANavy,left=13pt,right=13pt,top=12pt,bottom=12pt]
{\color{white}\sffamily\bfseries\large Ricostruzione DDTA del progetto DermaTriage}\par
\vspace{2mm}
{\color{white!88}\small La documentazione di progetto occupa la parte principale della pagina. La Base Analysis viene riportata in parallelo nella colonna laterale soltanto quando esiste un significato documentale da analizzare.}
\end{tcolorbox}

\vfill
\begin{ddtacore}[title={Struttura del case study}]
Il documento segue la progressione DDTA dal contesto generale del progetto agli elementi documentali successivi. Ogni elemento viene presentato nella colonna documentale; la colonna Base Analysis mostra esclusivamente l'analisi derivabile da quel significato, senza introdurre fatti di progetto aggiuntivi.
\end{ddtacore}

\vspace{5mm}
{\sffamily\scriptsize\color{DDTAGray}MD5 contenuto pagina 1: \texttt{\PageMDfive{1}}}\par
\end{titlepage}
%</PAGE-CONTENT:1>

%<PAGE-DESC:2>Contesto generale del progetto - Project Problem Framing
%<PAGE-CONTENT:2>
\PageHashFooter{\PageMDfive{2}}
\section{Contesto generale del progetto}

\setlength{\columnsep}{7mm}
\setlength{\columnseprule}{0.35pt}
\columnratio{0.67}
\begin{paracol}{2}

% LEFT 2/3 - DOCUMENTATION
{\sffamily\large\bfseries\color{DDTANavy}Documentazione DDTA}\par
\vspace{0.5em}

\begin{ddtacore}[title={Project Problem Framing}]
Il problema che DermaTriage deve affrontare e' il supporto al triage precoce di casi dermatologici relativi a lesioni potenzialmente oncologiche. Il progetto parte dalle informazioni gia' disponibili sul caso per determinarne l'urgenza e una priorita' operativa di presa in carico, con l'obiettivo di favorire un instradamento specialistico tempestivo.
\end{ddtacore}

\switchcolumn

% RIGHT 1/3 - BASE ANALYSIS ONLY
{\sffamily\large\bfseries\color{DDTABlue}Base Analysis}\par
\vspace{0.5em}

\begin{ddtaworking}[title={Base Analysis}]
\small
Il contesto generale definisce il problema affrontato dal progetto, ma non introduce ancora un elemento documentale strutturato da decomporre in referenti e proposizioni. La Base Analysis viene quindi applicata a partire dagli elementi DDTA successivi, quando il significato governato rende identificabili responsabilita', relazioni, condizioni o risultati analizzabili.
\end{ddtaworking}

\end{paracol}
%</PAGE-CONTENT:2>

\clearpage
%<PAGE-DESC:3>MR-01 - Valutazione di triage del caso dermatologico
%<PAGE-CONTENT:3>
\PageHashFooter{\PageMDfive{3}}
\section{MR-01 - Valutazione di triage del caso dermatologico}

\setlength{\columnsep}{7mm}
\setlength{\columnseprule}{0.35pt}
\columnratio{0.67}
\begin{paracol}{2}

{\sffamily\large\bfseries\color{DDTANavy}Documentazione DDTA}\par
\vspace{0.5em}

\begin{ddtacore}[title={MacroRequirement}]
\textbf{Intent}\par
Determinare, a partire dalle informazioni disponibili sul caso dermatologico, l'urgenza del caso e la relativa priorita' operativa di triage.

\textbf{Context}\par
La valutazione di triage parte dalle informazioni disponibili sul caso. Quando non e' presente un'immagine della lesione, il progetto prevede un percorso basato sui sintomi disponibili.

\textbf{Stakeholders}\par
Paziente.

\textbf{Scope}\par
\textbf{IN:} determinazione dell'urgenza del caso e della relativa priorita' operativa di triage, anche quando l'immagine non e' presente.\par
\textbf{OUT:} indicazione della destinazione specialistica; validazione o correzione medica dell'esito; adattamento successivo del comportamento del sistema sulla base della revisione clinica.

\textbf{Assumptions / Constraints}\par
--

\textbf{dependsOn}\par
\texttt{None}
\end{ddtacore}

\switchcolumn

{\sffamily\large\bfseries\color{DDTABlue}Base Analysis}\par
\vspace{0.5em}

\begin{ddtaworking}[title={Base Analysis}]
\small
La Base Analysis di questo MacroRequirement non e' ancora stata eseguita. Questa colonna resta intenzionalmente non popolata nel presente passaggio e verra' compilata nel successivo step di analisi, senza modificare il significato documentale dell'MR.
\end{ddtaworking}

\end{paracol}
%</PAGE-CONTENT:3>

\clearpage
%<PAGE-DESC:4>MR-02 - Indicazione della destinazione specialistica
%<PAGE-CONTENT:4>
\PageHashFooter{\PageMDfive{4}}
\section{MR-02 - Indicazione della destinazione specialistica}

\setlength{\columnsep}{7mm}
\setlength{\columnseprule}{0.35pt}
\columnratio{0.67}
\begin{paracol}{2}

{\sffamily\large\bfseries\color{DDTANavy}Documentazione DDTA}\par
\vspace{0.5em}

\begin{ddtacore}[title={MacroRequirement}]
\textbf{Intent}\par
Determinare un'indicazione di destinazione specialistica per il caso dermatologico, a supporto del successivo instradamento verso l'assistenza appropriata.

\textbf{Context}\par
DermaTriage associa al risultato del processo di triage un'informazione orientata alla destinazione specialistica. La documentazione disponibile non stabilisce il vocabolario delle destinazioni, la regola di selezione ne' la responsabilita' di assegnazione o prenotazione.

\textbf{Stakeholders}\par
--

\textbf{Scope}\par
\textbf{IN:} produzione di un'indicazione di destinazione specialistica per il caso.\par
\textbf{OUT:} prenotazione; assegnazione effettiva; tassonomia completa delle destinazioni; regola completa di selezione e fallback.

\textbf{Assumptions / Constraints}\par
--

\textbf{dependsOn}\par
\texttt{None}
\end{ddtacore}

\begin{ddtacheck}[title={STOP AT MR}]
Il significato macro e' sufficientemente distinguibile, ma la documentazione disponibile non governa commitment ulteriori sufficienti per una decomposizione fedele. Il branch si arresta intenzionalmente a MacroRequirement.
\end{ddtacheck}

\switchcolumn

{\sffamily\large\bfseries\color{DDTABlue}Base Analysis}\par
\vspace{0.5em}

\begin{ddtaworking}[title={Base Analysis}]
\small
La Base Analysis di questo MacroRequirement non e' ancora stata eseguita. Il successivo passaggio analitico dovra' consumare soltanto il significato qui documentato, preservando lo \texttt{STOP AT MR} e senza inferire regole di selezione specialistica mancanti.
\end{ddtaworking}

\end{paracol}
%</PAGE-CONTENT:4>

\clearpage
%<PAGE-DESC:5>MR-03 - Validazione e correzione medica degli esiti
%<PAGE-CONTENT:5>
\PageHashFooter{\PageMDfive{5}}
\section{MR-03 - Validazione e correzione medica degli esiti}

\setlength{\columnsep}{7mm}
\setlength{\columnseprule}{0.35pt}
\columnratio{0.67}
\begin{paracol}{2}

{\sffamily\large\bfseries\color{DDTANavy}Documentazione DDTA}\par
\vspace{0.5em}

\begin{ddtacore}[title={MacroRequirement}]
\textbf{Intent}\par
Consentire al medico di validare o correggere gli esiti prodotti da DermaTriage e rendere disponibile il risultato della revisione nel contesto del caso cui si riferisce.

\textbf{Context}\par
Gli esiti prodotti da DermaTriage possono essere sottoposti a revisione medica. Il medico puo' validarli o correggerli; il progetto gestisce il risultato della revisione senza includere in questa responsabilita' il successivo adattamento del sistema.

\textbf{Stakeholders}\par
Medico.

\textbf{Scope}\par
\textbf{IN:} gestione del risultato della validazione o correzione medica e sua associazione al caso o all'esito cui si riferisce.\par
\textbf{OUT:} meccanismi tecnici di scambio e autenticazione; adattamento successivo del comportamento del sistema sulla base delle correzioni raccolte.

\textbf{Assumptions / Constraints}\par
--

\textbf{dependsOn}\par
\texttt{None}
\end{ddtacore}

\switchcolumn

{\sffamily\large\bfseries\color{DDTABlue}Base Analysis}\par
\vspace{0.5em}

\begin{ddtaworking}[title={Base Analysis}]
\small
La Base Analysis di questo MacroRequirement non e' ancora stata eseguita. Nel passaggio successivo dovra' essere verificato quali referenti e relazioni siano realmente derivabili dal significato documentato, senza introdurre una semantica della revisione clinica non stabilita dalla documentazione.
\end{ddtaworking}

\end{paracol}
%</PAGE-CONTENT:5>

\clearpage
%<PAGE-DESC:6>MR-04 - Adattamento controllato sulla base della revisione clinica
%<PAGE-CONTENT:6>
\PageHashFooter{\PageMDfive{6}}
\section{MR-04 - Adattamento controllato sulla base della revisione clinica}

\setlength{\columnsep}{7mm}
\setlength{\columnseprule}{0.35pt}
\columnratio{0.67}
\begin{paracol}{2}

{\sffamily\large\bfseries\color{DDTANavy}Documentazione DDTA}\par
\vspace{0.5em}

\begin{ddtacore}[title={MacroRequirement}]
\textbf{Intent}\par
Consentire l'adattamento controllato del comportamento di DermaTriage utilizzando i risultati delle validazioni e correzioni mediche raccolte.

\textbf{Context}\par
Le correzioni mediche vengono utilizzate per modificare nel tempo il comportamento del sistema. L'evoluzione dei prompt e il riaddestramento del classificatore sono modalita' correnti di realizzazione di questa responsabilita' e non ne definiscono l'identita' macro.

\textbf{Stakeholders}\par
Medico.

\textbf{Scope}\par
\textbf{IN:} adattamento del comportamento del sistema sulla base dei risultati della revisione medica disponibili.\par
\textbf{OUT:} produzione del giudizio di validazione o correzione medica; dettagli dei singoli meccanismi di adattamento.

\textbf{Assumptions / Constraints}\par
--

\textbf{dependsOn}\par
\texttt{MR-03}
\end{ddtacore}

\switchcolumn

{\sffamily\large\bfseries\color{DDTABlue}Base Analysis}\par
\vspace{0.5em}

\begin{ddtaworking}[title={Base Analysis}]
\small
La Base Analysis di questo MacroRequirement non e' ancora stata eseguita. Il successivo passaggio dovra' preservare la dipendenza semantica da \texttt{MR-03} senza trasformarla in containment o co-ownership e senza promuovere automaticamente i meccanismi correnti di adattamento a responsabilita' macro autonome.
\end{ddtaworking}

\end{paracol}
%</PAGE-CONTENT:6>

\clearpage
%<PAGE-DESC:7>Indice di integrita' delle pagine
%<PAGE-CONTENT:7>
\PageHashFooter{\PageMDfive{7}}
\section{Indice di integrita' delle pagine}

\begin{ddtacore}[title={Regola di verifica}]
Ogni pagina possiede un MD5 del proprio contenuto canonico. Il calcolo esclude il valore MD5 stampato sulla pagina, gli header/footer e i metadati di impaginazione esterni al relativo blocco di contenuto. Per la pagina dell'indice, gli hash delle pagine precedenti fanno parte del contenuto controllato; il solo hash della pagina indice viene sostituito dal token canonico \texttt{<SELF>} durante il calcolo, evitando una dipendenza circolare.
\end{ddtacore}

\vspace{3mm}
\begin{tabularx}{\textwidth}{L{12mm}L{72mm}Y}
\toprule
\textbf{Pag.} & \textbf{Contenuto} & \textbf{MD5 contenuto} \\
\midrule
%<PAGE-INTEGRITY-ROWS>
\IntegrityRow{1}{Frontespizio e struttura del case study}{\PageMDfive{1}}
\IntegrityRow{2}{Contesto generale del progetto - Project Problem Framing}{\PageMDfive{2}}
\IntegrityRow{3}{MR-01 - Valutazione di triage del caso dermatologico}{\PageMDfive{3}}
\IntegrityRow{4}{MR-02 - Indicazione della destinazione specialistica}{\PageMDfive{4}}
\IntegrityRow{5}{MR-03 - Validazione e correzione medica degli esiti}{\PageMDfive{5}}
\IntegrityRow{6}{MR-04 - Adattamento controllato sulla base della revisione clinica}{\PageMDfive{6}}
\IntegrityRow{7}{Indice di integrita' delle pagine}{\PageMDfive{7}}
%</PAGE-INTEGRITY-ROWS>
\bottomrule
\end{tabularx}

\vspace{4mm}
{\small\color{DDTAGray}
Quando una pagina viene modificata intenzionalmente, cambia il suo MD5. Le pagine non interessate dal passaggio devono mantenere lo stesso valore; una variazione inattesa segnala una modifica da riesaminare prima di proseguire.}
%</PAGE-CONTENT:7>

\end{document}
